% ============================================================================
%  GENERATED FILE — do not edit.
%
%  Produced by scripts/96_flatten_paper.py from main.tex and sections/*.tex.
%  This is main.tex with every \input inlined, so the whole manuscript can be
%  read or compiled as a single document (e.g. pasted into Overleaf).
%
%  To change the text, edit the section file named in the marker comment above
%  the passage, then re-run:
%
%      python scripts/96_flatten_paper.py
%
%  Editing this file directly will be overwritten.
% ============================================================================

% ============================================================================
%  When compliance cannot be determined
%  Draft v0.1 --- 2026-08-03
%
%  BUILD:  latexmk -pdf main.tex
%
%  RULE FOR THIS MANUSCRIPT
%  ------------------------
%  Every number in the text is produced by a script in ../scripts/ and appears
%  in ../eval/. Nothing is typed from memory. Values still to be computed are
%  marked \pending{...} and MUST NOT survive to submission; \pending renders in
%  red so they cannot be missed.
%
%  Provenance of each figure is given in its caption.
% ============================================================================

\documentclass[11pt,a4paper]{article}

\usepackage[utf8]{inputenc}
\usepackage[T1]{fontenc}
\usepackage[english]{babel}
\usepackage{lmodern}
\usepackage{microtype}
\usepackage[margin=2.5cm]{geometry}
\usepackage{graphicx}
\graphicspath{{figures/}}
\usepackage{booktabs}
\usepackage{longtable}
\usepackage{amsmath,amssymb}
\usepackage{siunitx}
\usepackage{xcolor}
\usepackage[hidelinks]{hyperref}
\usepackage{authblk}
\usepackage{lineno}

\sisetup{detect-weight=true, per-mode=symbol}

% Unresolved values. Deliberately loud.
\newcommand{\pending}[1]{\textcolor{red}{\textbf{[TO COMPUTE: #1]}}}
% Values that came from a script; the argument is the file that produced them.
\newcommand{\src}[1]{\,\footnote{Produced by \texttt{#1}.}}

% ---------------------------------------------------------------------------
% TITLE --- settled. "Not detected is not absent" states the thesis in four words
% and is legible to a chemist and a computer scientist on first reading.
% "Micropollutant" is the concrete, searchable term; "compliance" was dropped as
% bureaucratic and narrow.
% ---------------------------------------------------------------------------
\title{\textbf{Not detected is not absent:\\
an ontology for detection limits and compliance\\
that cannot be determined in micropollutant monitoring}}

\author[1]{Recep Can Altınbağ}
\author[2]{Mehmet Tahir Sandıkkaya}
\author[1]{Ulaş Pekyavaş}
\affil[1]{Institute of Environmental Sciences, Bo\u{g}azi\c{c}i University,
TR34342 Be\c{s}ikta\c{s}, Istanbul, T\"urkiye}
\affil[2]{Department of Computer Engineering, Istanbul Technical University,
TR34469 Sar{\i}yer, Istanbul, T\"urkiye}

\date{Draft --- \today}

\begin{document}
\maketitle
\linenumbers

\begin{abstract}
% >>> begin sections/00-abstract.tex
A large share of micropollutant measurements are \emph{non-detections}: the
analyte was sought and not found. Such a result establishes an upper bound, not
a value, yet records store it as a number, or without the limit that bounds it.
European law already requires the distinction --- Article~3(3b) of Directive
2008/105/EC excludes a below-quantification result from status assessment where
that limit exceeds the standard --- but a two-valued model cannot express it, so
compliance with the provision cannot be audited from the record as kept.

Auditing the European Environment Agency's Waterbase (\num{4190833} river
station-years, \num{637} substances, \num{37} countries) shows
\num{44.9}\% of samples below the quantification limit, \num{25.5}\% of
station-years recording neither a censoring flag nor a limit, and
\num{46.6}\% of assessments made with a method failing the legal performance
criterion of Directive 2009/90/EC. The cost is not only hidden failure. Over
\num{40000} assessments, substituting a number for a non-detection yields
between \num{3382} and \num{12042} exceedances depending only on whether the
substitution is at zero, half the limit or the limit, while just \num{908} rest
on a quantified measurement: the enforcement-relevant count moves by a factor of
\num{3.6} with a rule the data does not fix.

We present CENSO, an OWL~2 ontology on SOSA/SSN whose single commitment is that
\emph{the quantification limit belongs to the analytical run, not the
instrument}. A non-detection is then typed as the interval it establishes, and
compliance becomes four-valued, adding \emph{possibly exceeding} and
\emph{cannot be determined}. Regulations are pluggable packages: applying two
jurisdictions to the same observations changes \num{18.5}\% of the outcomes both
regulate, with no code change. Ontology, packages, pipeline and per-figure
data are released openly.
% <<< end sections/00-abstract.tex
\end{abstract}

\noindent\textbf{Keywords:} censored data; limit of detection; environmental
quality standards; ontology; open world assumption; micropollutants;
compliance assessment; river basin monitoring

\begin{figure}[htbp]\centering
\includegraphics[width=\textwidth]{fig01_graphical_abstract.pdf}
\caption{Graphical abstract. Five heterogeneous inputs are expressed in one
knowledge graph, from which reach-level verdicts are derived. The fourth
verdict, \emph{cannot be determined}, is the one no conventional pipeline
can return.
Generated by \texttt{scripts/90\_figures.py}.}
\label{fig:abstract}
\end{figure}

% >>> begin sections/01-introduction.tex
\section{Introduction}

Trace organic contaminants -- pharmaceuticals, pesticides, biocides, UV filters,
benzotriazoles, flame retardants -- are ubiquitous in surface waters. Two
decades after Schwarzenbach et al.~\cite{schwarzenbach2006challenge} set out the
challenge they pose, the scale of the problem is established rather than
suspected: a survey of \num{1052} sites on \num{258} rivers across \num{104}
countries found pharmaceutical contamination on every continent, exceeding
safe concentrations at a quarter of them~\cite{wilkinson2022pharma}, and an
EU-wide screening detected polar organic pollutants in essentially every river
sampled~\cite{loos2009eu}. The effects are not confined to chemistry. Pesticide
concentrations at levels European regulation treats as acceptable are
associated with losses of regional stream invertebrate
biodiversity~\cite{beketov2013biodiversity}, and the substances that reach
surface waters continue to outpace the frameworks built to manage
them~\cite{geissen2015emerging}.

These substances are regulated in the European Union through environmental
quality standards (EQS) under the Water Framework Directive~\cite{eu2000wfd}, transposed in
T\"urkiye as \emph{\c{C}evresel Kalite Standartlar{\i}} (\c{C}KS). Compliance is
assessed by comparing measured concentrations against those standards -- and
that comparison is where this paper begins, because it is less well defined
than it appears.

Most such measurements are non-detections. The analyte is sought and not found,
which establishes an \emph{upper bound} on its concentration and nothing more.
Crucially, it does not establish absence: the substance may be present below the
limit of detection (LOD) of the method. The statement is open-world.

Reporting practice replaces that statement with a number, usually zero or half
the limit of quantification (LOQ). Helsel showed two decades ago that
substitution fabricates data and distorts every derived statistic
\cite{helsel2006fabricating}, and the point is not contested. It is nonetheless
still standard, and this paper begins from the observation that the reason is
structural rather than merely habitual: \textbf{the data models in use cannot
hold the statement}. A relational table operates under the closed-world
assumption, in which an absent assertion is false; a column of concentrations
has no cell for ``present or absent, below \num{0.02}\,\si{\micro\gram\per\litre}''.
Substituting a zero is therefore not a rounding convention but a type error,
and what it destroys is the modality of the measurement, not its magnitude.

The consequence is not confined to statistics. Regulatory assessment implicitly
requires three verdicts -- compliant, exceeding, and \emph{cannot be determined}
-- while the systems that produce it can express only the first two. The third
case is not hypothetical. Where the LOQ of a method lies above the EQS of the
substance it measures, a laboratory can correctly report ``not detected'' for a
sample that may exceed the legal limit by two orders of magnitude, and a
closed-world pipeline will record the result as compliance.

This paper makes that cost measurable at the scale the law operates on. We
analyse the European Environment Agency's Waterbase, the compilation of
surface-water monitoring that Member States report under the
Directive~\cite{eea_waterbase}: \num{4190833} river station-years over
\num{637} substances and \num{37} reporting countries, representing
\num{32993382} samples. Waterbase records, for each station, substance and year,
both a below-quantification flag and the limit that was in force --- the two
fields whose absence is the subject of this paper --- so where they are missing,
that absence is itself the measurement.

Our contributions are:

\begin{enumerate}
  \item \textbf{An implementability gap in the law.} Article~3(3b) of the
    consolidated Directive requires that a below-quantification result whose
    limit exceeds the standard ``shall not be considered'' for status
    assessment, and Article~4(1) of Directive 2009/90/EC requires that limit to
    sit at or below \num{30}\,\% of the standard. Both mandate an outcome that a
    two-valued data model cannot represent, so neither can be audited from the
    data as recorded.
  \item \textbf{A pan-European audit.} We quantify how much of the reported
    record cannot support the conclusions ordinarily drawn from it: the share of
    samples below the quantification limit, the share of station-years carrying
    neither a censoring flag nor a limit, and the share assessed with a method
    that fails the legal performance criterion.
  \item \textbf{An explanation.} We show that the gap in existing vocabularies
    follows from a measurement model: SOSA/SSN and its descendants were designed
    for in-situ \emph{sensing}, in which a detection limit is a property of the
    device. Trace organic monitoring is \emph{sampling followed by laboratory
    analysis}, in which the limit belongs to the analytical run. The field
    inherited the first model and did not adapt it.
  \item \textbf{An instrument.} We present CENSO, an ontology that binds the
    limit to the analytical run, records why a result is an interval, treats a
    threshold as a conditional judgement with preconditions, and admits an
    outcome for a verdict that cannot be determined. Its axioms are tested
    against defects drawn from the
    data, and it is exercised on the public record rather than on a private
    survey.
  \item \textbf{A comparison.} We assess \num{18} published ontologies by
    parsing each one, rather than by reading its documentation, and report the
    method and its corrections so the table can be re-derived.
\end{enumerate}

The ontology is the instrument, not the result. Without a vocabulary that
separates \emph{compliant} from \emph{cannot be determined}, none of the
quantities reported here can be counted at all.
% <<< end sections/01-introduction.tex
% >>> begin sections/02-related-work.tex
\section{Related work}

\subsection{Censored environmental data}
The statistical treatment of left-censored measurements is long settled.
Substituting a constant for a non-detect biases means, variances and trends, and
the size of the bias depends on the censoring rate, which varies between
analytes and campaigns \cite{helsel2006fabricating,helsel2012nondetects}.
Robust alternatives -- regression on order statistics, maximum likelihood,
survival methods -- are available and unused in routine practice. Our concern is
upstream of the statistics: by the time an analyst opens the table, the
censoring status has usually already been discarded, and no method can recover
what the schema could not hold.

\subsection{Observation vocabularies}
SOSA/SSN is the W3C/OGC standard for sensors and observations
\cite{janowicz2019sosa,haller2019ssn}. \texttt{sosa:Observation} carries
\texttt{hasSimpleResult}, a point value. The SSN Systems module does define
\texttt{ssn-system:DetectionLimit}, but as a subclass of
\texttt{ssn-system:SystemProperty}: it describes what a \emph{sensor} can do.
It is never attached to an observation or its result, so nothing in the standard
can state that a particular measurement fell below it.

That is coherent for the modality SOSA was designed for. A sensor observes in
situ and continuously, and its detection limit really is a fixed property of the
device. Trace organic monitoring is a different modality: a physical sample is
taken at an instant, transported, and analysed later by, for example, LC--MS/MS,
in a run with its own calibration curve and hence its own limits for that
analyte and matrix. SOSA supplies \texttt{Sampling}, \texttt{Sampler} and
\texttt{Sample}, but nothing links the analysis of a sample to a detection
limit. We argue in Section~\ref{sec:results-gap} that the absence throughout the
water-ontology literature follows from this inherited model.

\subsection{Water domain ontologies}
Several ontologies address water quality: WaWO+ for reasoning about river basins
\cite{olivafelipe2017wawoplus}, the observational process ontology OPO
\cite{wang2020opo}, the InWaterSense expert system
\cite{jajaga2015expert,ahmedi2013wqm}, SAREF4WATR for water-domain IoT
\cite{etsi2025saref4watr}, and the Water Health Open Knowledge Graph
\cite{whow2023onto}. Query technologies over SSN have also been developed
\cite{calbimonte2012ssw}.

WHOW is the closest prior work and, unlike most, publishes its ontology at a
dereferenceable IRI. It defines a \texttt{Range} class with \texttt{lowerBound}
and \texttt{upperBound} attached to \texttt{ObservationValue}, and its own
documentation gives the example of a copper concentration reported as
\texttt{<10}, meaning $[0,10]$. \textbf{WHOW therefore does represent a censored
result as an interval}, and this work extends it rather than filling a void.
What WHOW does not record is why the interval exists: nothing distinguishes
censoring from measurement spread, nothing ties the interval to the limit that
produced it, and the ontology has no compliance vocabulary, so an
indeterminate
outcome cannot be posed.

Section~\ref{sec:results-gap} reports a systematic comparison of \num{18}
ontologies, each parsed from its published file.

\paragraph{Analytical coverage already changes conclusions.}
That the reach of a method decides what a survey can conclude is established.
Moschet et al.~\cite{moschet2014screening} screened Swiss streams for
\num{249} pesticides and found that the substances a routine programme omits
dominate the risk picture: broadening the analytical window changed the quality
assessment outright. Malaj et al.~\cite{malaj2014continental} used the same
Waterbase records analysed in Section~\ref{sec:results-waterbase} to show that
organic chemicals threaten freshwater ecosystems at continental scale. Kase et
al.~\cite{kase2018estrogens} state the problem in its sharpest form for
steroidal estrogens: for several the required standard lies below what routine
laboratories can quantify at all, so compliance monitoring must either improve
the method or accept that it cannot answer the question.

Our contribution is adjacent rather than competing. Those studies ask what
better chemistry would reveal; we ask what the record, as it stands, is capable
of establishing -- and show that the answer is often nothing, in a way no
substitution rule repairs. Ahkola et al.~\cite{ahkola2024uncertainty} survey
uncertainty in micropollutant modelling and treat it as a quantity to be
estimated and propagated. The states we identify are not uncertain in that
sense: where the threshold falls inside $[0,\mathrm{LOD}]$ there is no
distribution to place over the verdict, because the measurement did not
discriminate between the outcomes. That is a representational gap, not an
estimation problem.

\paragraph{Semantic infrastructure for water.}
Systematic review of the domain's vocabularies~\cite{water_onto_review} finds
no shortage of them; the Water Health Open Knowledge Graph~\cite{whokg2025} and
recent interoperability work in water treatment~\cite{saka2026acquirium} show
the direction of travel. What none of them carries is the epistemic status of
the measurement itself, which is the gap Section~\ref{sec:results-gap}
quantifies by parsing the files.

\paragraph{A fourth reason the Directive misses risk.}
Weisner et al.~\cite{weisner2022threereasons} give three reasons the Water
Framework Directive fails to identify pesticide risk: the list of analytes
monitored is outdated and inconsistent, the choice of sampling sites is biased,
and relevant thresholds outside the Directive go unconsidered. Each concerns
\emph{which} substances are looked for and \emph{where}. We report a fourth,
which applies to the substances that are on the list and the sites that are
monitored: even there, the record frequently cannot decide the question it was
collected to answer. That the analytical capability required by the Directive
was uncertain from the outset is not a new observation --
Coquery et al.~\cite{coquery2005analytical} raised it when the priority
substance list was first being implemented, and the Commission's own review of
Watch List methods restates it for the current
substances~\cite{loos2024watchlist}. What has been missing is a way to
\emph{count} the consequence, which requires representing the failure rather
than substituting past it.
% <<< end sections/02-related-work.tex
% >>> begin sections/03-materials.tex
\section{Data and regulatory sources}
\label{sec:materials}

\subsection{EEA Waterbase}
The measurements are the European Environment Agency's Waterbase --- Water
Quality ICM, aggregated release, which compiles the surface-water monitoring
that Member States report under the Water Framework
Directive~\cite{eea_waterbase,eu2000wfd}. One row is one station, substance and
year, carrying the annual mean, the number of samples behind it, a flag for
whether that mean fell below the limit of quantification, and
\texttt{procedureLOQValue}, the limit in force. Those last two fields are what
make this the right dataset to test on: they are exactly the information whose
absence this paper is about, so where they are missing that is itself the
observation.

River water only is retained. An inland-surface-water standard does not apply to
lakes, transitional or coastal waters, and comparing across them would be a
category error rather than a finding. That leaves \num{4190833} station-years
over \num{637} substances and \num{37} reporting countries, representing
\num{32993382} underlying samples. Station names and coordinates come from the
companion \texttt{SpatialObjects} release.

The aggregated release is used rather than the disaggregated one for a reason of
kind, not of size. An annual-average environmental quality standard is defined
against an annual mean, which is precisely what one aggregated row holds.
Counting individual samples against such a standard would compare two quantities
the regulation never intended to be compared.

\subsection{Regulatory thresholds}
Every threshold is transcribed from its primary text, never from a secondary
compilation.

\paragraph{Environmental quality standards.}
Annex~I of Directive 2008/105/EC in the consolidated version in force on 10 May
2026~\cite{eu2008eqs,eu2013priority}, which incorporates the 2026 update adding
PFAS, pharmaceuticals, bisphenols and pesticide metabolites. Parsing it is not
mechanical: the table wraps one logical row over many physical lines, entries
carry several CAS numbers, and scientific-notation values are rendered with
spaces inside the mantissa, so that ``$6 \times 10^{-4}$'' tokenises as four
separate numbers\src{scripts/10\_parse\_eu\_eqs.py}. Group standards --- which
apply to a sum rather than to any single substance --- are flagged and excluded
from per-substance comparison; applying the total-pesticide limit to each of the
\num{75} substances it names would compare one substance against a limit written
for their aggregate.

\paragraph{Analytical performance.}
Directive 2009/90/EC sets what a method must achieve before its results may be
used at all~\cite{eu2009qaqc}. Article~4(1) requires ``a limit of quantification
equal or below a value of \num{30}\,\% of the relevant environmental quality
standards''. That is a stricter test than asking whether the limit exceeds the
standard outright, and Section~\ref{sec:results-waterbase} reports both.

\paragraph{A second package, for swappability.}
Table~4 of the Turkish \emph{Yer\"ust\"u Su Kalitesi Y\"onetmeli\u{g}i} as
amended in \num{2016}~\cite{tr2012ysky} is transcribed by the same
route\src{scripts/08\_verify\_thresholds.py} and released alongside. It is not
analysed here. It exists so that the claim in
Section~\ref{sec:methods-packages} --- that a regulation is a module which can
be swapped --- rests on two jurisdictions actually being loadable rather than on
an assurance that they would be.

\subsection{What is joined to what}
Substances are matched to thresholds by CAS registry number, which Waterbase
carries in \texttt{observedPropertyDeterminandCode}. Name matching was not used:
chemical names diverge between a directive, a national register and a
laboratory's reporting sheet, and a missed join would silently understate every
count that follows. Concentrations arrive in one column as
\si{\milli\gram\per\litre}, \si{\micro\gram\per\litre} and
\si{\nano\gram\per\litre}; each is converted before any comparison, since
treating them as interchangeable would be wrong by three orders of magnitude in
either direction.
% <<< end sections/03-materials.tex
% >>> begin sections/04-methods.tex
\section{Methods}
\label{sec:methods}

\subsection{CENSO: binding the limit to the result}
CENSO\footnote{\url{https://w3id.org/censo/}, CC~BY~4.0.} extends
SOSA/SSN~\cite{janowicz2019sosa,haller2019ssn} for observations produced by a
method that has a detection limit. It reuses PROV-O~\cite{w3c2013provo} for the
analytical run and the provenance of a transcribed threshold,
SKOS~\cite{w3c2009skos} for substance concepts and QUDT for units; substances
are aligned to ChEBI~\cite{degtyarenko2008chebi} where an identifier exists
(\num{180} of \num{240}). It was developed following the
LOT~\cite{poveda2022lot} methodology, a lighter-weight successor to NeOn~\cite{suarez2012neon}, whose competency-question-first sequence
is what produced the evaluation in Section~\ref{sec:results-cq}, and it is
published to be findable and reusable in the sense
of~\cite{wilkinson2016fair}. Four commitments
distinguish it from the vocabularies surveyed in
Section~\ref{sec:results-gap}.

\paragraph{The limit belongs to the analytical run.}
\texttt{censo:AnalyticalMethod} is a \texttt{sosa:Procedure} carrying
\texttt{limitOfDetection}, \texttt{limitOfQuantification} and a unit;
\texttt{censo:AnalyticalRun} reifies one execution of it, with its own
calibration curve and therefore its own limits, and
\texttt{censo:MassSpectrometryTransition} records the precursor and product
ions and retention time that identify the measurement channel. An observation
links to the run that produced it, so the limits in force for that measurement
are recoverable. This is the level at which analytical provenance exists in
practice and it is absent from every ontology we surveyed.

\paragraph{Detection status is explicit and exhaustive.}
\texttt{CensoredObservation}, \texttt{EstimatedObservation},
\texttt{QuantifiedObservation} and \texttt{UnresolvedObservation} are pairwise
disjoint and cover \texttt{AssessedObservation}. The absence of an observation
individual is a fifth, distinct situation -- never measured -- and is
deliberately not modelled as a class: under the open-world assumption, not
asserting an observation is the correct way to say it was not made. This is the
distinction a \texttt{WHERE concentration = 0} query cannot draw, and
conflating it manufactures findings.

\paragraph{A threshold is a conditional judgement.}
\texttt{censo:Threshold} carries a value, a unit, a citation to the regulation,
and \texttt{requiresCondition} links to \texttt{ApplicabilityCondition}
subclasses: matrix, analytical fraction, bioavailability basis, hardness class.
Where a precondition is unmet the outcome is \texttt{PreconditionUnmet}, a
subclass of \texttt{IndeterminateCompliance} -- not a default pass.

\paragraph{Compliance is four-valued, per observation and threshold.}
\texttt{Exceedance}, \texttt{PossibleExceedance}, \texttt{Compliant} and
\texttt{IndeterminateCompliance} are mutually exclusive and jointly exhaustive
\emph{for a given observation--threshold pair}. The qualification is not
pedantry and it is the reason no class-level disjointness axiom is asserted: one
measurement compared against two jurisdictions' standards may legitimately be
compliant under one and indeterminate under the other, and
Section~\ref{sec:results-dual} depends on the graph being able to hold both at
once. Exclusivity is therefore carried by an \texttt{owl:AllDisjointProperties}
axiom over \texttt{exceeds}, \texttt{possiblyExceeds} and
\texttt{belowThreshold}, which forbids two verdicts about the same pair while
permitting different verdicts about different thresholds. Exhaustiveness is a
closed-world statement and is materialised in SHACL, not asserted in OWL, for
the reason given below.

The first three are defined classes over materialised comparison properties;
the fourth has four documented causes: no threshold exists, the method's LOQ
exceeds the threshold, a precondition is unmet, or the observation is censored
and the threshold lies inside $[0, \mathrm{LOD}]$.

Figure~\ref{fig:geometry} is the whole argument in one picture. A censored
result is not a point on the concentration axis but the interval
$[0,\mathrm{LOD}]$, and a threshold is a line that either clears that interval
or falls inside it. Where it clears (a), the non-detection \emph{does} decide
the question, and the ontology returns \texttt{Compliant} -- censoring is not
an excuse for refusing every verdict. Where it falls inside (b), no verdict
exists to be found: the measurement is consistent with both compliance and
exceedance, and any substitution rule merely picks one. Where the method's
quantification limit lies above the threshold (d), no result the method could
have returned would have decided it, so the defect is analytical rather than
representational and no reprocessing repairs it. A two-valued model must assign
(b) and (d) to \texttt{Compliant}, which is why the loss is invisible in
current practice.

\begin{figure}[htbp]\centering
\includegraphics[width=\textwidth]{fig06_decision_geometry.pdf}
\caption{Why two verdicts are not enough. A non-detection is the interval
$[0,\mathrm{LOD}]$ (shaded), not a value; the threshold $T$ is a line. Panels
(a) and (c) are decidable and a conventional pipeline gets them right. Panels
(b) and (d) are not, and a conventional pipeline silently records them as
compliant. The distinction is a property of where $T$ falls relative to the
method's limits, so it is available to a reasoner but invisible to a numeric
column. Generated by \texttt{scripts/90\_figures.py}.}
\label{fig:geometry}
\end{figure}

\subsection{Regulations as pluggable, versioned packages}
\label{sec:methods-packages}
Thresholds expire. \texttt{cereg:RegulationPackage} makes each regulation a
self-contained module contributing only individuals, with jurisdiction, dates,
a citation and a dereferenceable source document, plus a machine-readable
\texttt{transcriptionStatus} so that unverified values can be excluded from
analysis rather than trusted. Group thresholds (\texttt{cereg:GroupThreshold})
carry an aggregation function and a \texttt{requiresCompleteGroup} flag, so a
sum over a substance group that the survey did not fully measure yields
\texttt{NotCoveredBySurvey} rather than an optimistic partial sum.

Two packages are released with the ontology\src{scripts/19\_build\_regulation\_packages.py}:
\texttt{tr-ysky-2016} (\num{217} analytes, \num{434} thresholds) and
\texttt{eu-2008-105-2026} (\num{62} analytes, \num{114} thresholds), both
transcribed from the primary legal texts and carrying
\texttt{VerifiedAgainstPrimarySource}. Each imports the vocabulary and
contributes \emph{individuals only} -- no class, no property -- which is what
makes one safe to load in place of, or alongside, another; the validator checks
this rather than assuming it. Because thresholds from different packages carry
different IRIs and different \texttt{definedBy} links, the Turkish and European
assessments of the same observations coexist in one graph instead of
overwriting each other.

\subsection{Is this not simply what detection limits are for?}
\label{sec:methods-noise}
The objection to answer first is that limits of detection exist precisely to
separate signal from noise, so a laboratory reporting ``not detected'' is doing
exactly the right thing and there is nothing here to fix. The analytical
practice is indeed correct, and nothing in this paper disputes it. But two
statements are being conflated:

\begin{itemize}
  \item \emph{the measurement cannot resolve a value} -- true below the limit,
    and the reason a number must not be reported; and
  \item \emph{the concentration is zero} -- not established by the measurement
    at all.
\end{itemize}

The laboratory asserts the first. The database records the second. What is lost
is not a value -- there was none to keep -- but the \emph{bound}: a
non-detection establishes $c \in [0,\mathrm{LOD}]$, which is a genuine
constraint, and \num{0.0} is not a weaker version of it but a different claim.

Whether that loss matters is decided by something the analyst never saw. If
$\mathrm{LOD} < T$ the non-detection settles the question and the verdict is
\emph{compliant}; noise filtering has cost nothing. If $\mathrm{LOD} > T$ the
same record settles nothing, because a concentration just below the limit may
lie well above the standard. The method was chosen on analytical grounds, and
the threshold arrives afterwards from a legal instrument; their relative
position is therefore accidental, and it is what determines whether ``not
detected'' is informative.

The regulation itself makes exactly this distinction, which is the strongest
evidence that the third outcome is not our invention. Article~3(3b) of the
consolidated Directive 2008/105/EC provides that where a result is reported as
less than the limit of quantification and \emph{that limit is above the EQS},
``the result for the substance being measured shall not be considered for the
purposes of assessing the overall chemical status of that water
body''~\cite{eu2008eqs,eu2009qaqc}. That is neither \emph{compliant} nor
\emph{exceeding}: the law requires a third outcome and requires it to be
identifiable. A two-valued schema cannot represent it, so the requirement
cannot be honoured, audited, or even counted.
Section~\ref{sec:results-waterbase} reports how often it arises.

\subsection{Could a column have done this?}
\label{sec:methods-why}
The objection must be met directly, because it is correct as far as it goes: a
relational table with a \texttt{censored} flag and an \texttt{lod} column
yields every number in Section~\ref{sec:results-waterbase}. Our own analysis scripts
are plain Python over CSV. We therefore do not claim that the quantities
reported here are otherwise uncomputable.

What the flag does not give is the following, and these are the grounds on which
the ontology earns its place.

\begin{enumerate}
  \item \textbf{Several regulations over one set of observations.} The Turkish
    and European assessments of these measurements differ substantially
    (Section~\ref{sec:results-waterbase}).
    With regulations as packages this is a re-query; with a threshold column it
    is a re-implementation, and the two results cannot be held simultaneously.
  \item \textbf{Preconditions belong to the threshold, not the measurement.}
    Hardness class, bioavailability covariates and analytical fraction condition
    whether a standard may be applied at all. A flag on the observation cannot
    express a condition on the standard.
  \item \textbf{Justification.} A defined class yields a derivation for each
    verdict, which is what a regulator can be shown. A boolean yields a value.
  \item \textbf{Transfer.} The pattern is domain-independent; a column in one
    team's spreadsheet is not.
\end{enumerate}

Of these, this paper demonstrates (1) and (2) on real data and (3) on the axiom
suite. Claim (4) is design intent, evidenced only by the vocabulary's
independence from the water domain, and we mark it as such.

\subsection{What OWL does, and what it does not}
Stated plainly because the distinction is often blurred. OWL supplies identity,
taxonomy, disjointness, classification and explanation. It cannot compare two
data values, so \texttt{exceeds}, \texttt{possiblyExceeds} and
\texttt{belowThreshold} are \emph{materialised} by the rule layer and then
classified by the ontology. And because the open-world assumption that makes
CENSO right for censoring also prevents ``none of the above'' being inferred,
\texttt{IndeterminateCompliance} is materialised under a locally closed world
using SHACL. We do not claim that a description-logic reasoner computes
compliance.

\subsection{Reproducibility}
Every number in this paper is produced by a script and written to a file under
\texttt{eval/}; none is typed by hand. The pipeline is deterministic, uses only
paths relative to the repository root, and runs in pure Python -- OWL~2~RL
entailment via \texttt{owlrl} and closed-world validation via
\texttt{pyshacl} -- so no JVM or external reasoner is required. External API
responses (PubChem, CrossRef) and every comparison ontology are cached in the
repository, so results can be re-derived offline after those services change.
The ontology was scanned with OOPS!~\cite{poveda2014oops}, and FOOPS!\
\cite{garijo2021foops} is the FAIR assessment we report as not yet run rather
than assert. A generated manifest, \texttt{eval/numbers\_manifest.md}, maps
every numeric
claim in this paper to the artefact and script that produced it, and fails if any
claim cannot be traced\src{scripts/95\_numbers\_manifest.py}.
% <<< end sections/04-methods.tex
% >>> begin sections/05-results.tex
\section{Results}

\subsection{How often the record cannot support its own verdict}
\label{sec:results-waterbase}
A representational claim is only worth as much as the evidence that the problem
it addresses is real and widespread, and evidence drawn from one laboratory's
reporting habits would be a finding about that laboratory. The EEA's Waterbase
records, for each station, substance and year, a below-quantification flag and
the limit of quantification in force -- the two fields whose absence is our
finding -- so the failures can be counted across an entire
continent\src{scripts/22\_waterbase\_external.py}. We retain river water only, since
an inland-surface-water standard does not apply to lakes or coastal waters, join
substances to the verified European thresholds by CAS number, and convert every
concentration to \si{\micro\gram\per\litre} before comparison. This leaves
\num{4190833} station-years over \num{637} substances and \num{37} reporting
countries, representing \num{32993382} underlying samples.

\paragraph{Censoring is the normal case, not an edge case.}
\textbf{\num{44.9}\% of all \num{32993382} samples lie below the limit of
quantification}, and \num{60.4}\% of station-years contain at least one such
sample. Censoring is not an edge case to be handled once the real analysis is
done; it is the ordinary condition of trace monitoring, and any treatment of it
propagates into most of the record. Where the record keeps its bound is not
randomly distributed: Figure~\ref{fig:europe-map} maps every station in the
graph, and the pattern follows national borders rather than river basins.

\begin{figure}[htbp]\centering
\includegraphics[width=0.86\textwidth]{fig02_europe_map.pdf}
\caption{The \num{7652} monitoring stations in the graph that carry
coordinates, coloured by the share of their reporting authority's station-years
that record no quantification limit. The pattern follows national borders
rather than river basins, which is what places the defect in reporting practice
rather than in chemistry or hydrology. Generated by
\texttt{scripts/90\_figures.py}.}
\label{fig:europe-map}
\end{figure}

\paragraph{The bound is missing for a quarter of the record.}
Where a reporter declares censoring, the limit is almost always recorded: only
\num{61034} of \num{3118074} declared station-years lack it. The failure is
elsewhere. \textbf{\num{1066797} station-years (\num{25.5}\%) declare nothing at
all} -- no flag and no limit -- so a zero and a measured trace are
indistinguishable, and no substitution rule can even be chosen because there is
nothing to substitute from. These are
\texttt{UnresolvedObservation} at continental scale. The EEA's own quality
control identifies the same defect independently, marking \num{330973} rows
\texttt{QC\_LOQ\_UNKNOWN}.

\begin{figure}[htbp]\centering
\includegraphics[width=0.86\textwidth]{fig03_country_practice.pdf}
\caption{Each point is a reporting country with at least \num{5000}
station-years; area is proportional to that count. The relationship runs the
wrong way for a chemical explanation: authorities recording least (right) are
found in breach least (bottom), because a limit that was never reported cannot
be tested against the standard. The measured failure rate is therefore a lower
bound. Generated by \texttt{scripts/90\_figures.py}.}
\label{fig:country-practice}
\end{figure}

\paragraph{Half of European monitoring fails the legal performance criterion.}
Article~4(1) of Directive 2009/90/EC sets the standard a method must meet before
its results may be used: ``a limit of quantification equal or below a value of
\num{30}\,\% of the relevant environmental quality
standards''~\cite{eu2009qaqc}. Measured against that criterion rather than
against the standard itself, \textbf{\num{46.6}\% of station-years assessed
against a European standard were produced by a method that does not satisfy
it}. This is not a defect of the result but of the monitoring: the requirement
applies to the method, before any sample is taken.

\paragraph{Article~3(3b) is routinely unenforceable.}
Of the station-years for which a European standard exists, \textbf{\num{23.6}\%
report a quantification limit above that standard}: \num{178963} assessments
that the Directive requires to be set aside and that a two-valued pipeline
records as compliant. For four pyrethroid insecticides the figure is
\num{100}\% -- bifenthrin, deltamethrin, esfenvalerate and cypermethrin have
\emph{never}, in this record, been measured by a method capable of deciding
their own standard. Chlorpyrifos reaches \num{91.9}\% over \num{25960}
station-years, and imidacloprid \num{75.8}\% --- an entire, newly regulated
insecticide class that European monitoring, as reported, cannot assess against
its own standard. Figure~\ref{fig:loqeqs} ranks the substances for which this
holds most often.

These are not the failures of one laboratory or one programme. They are
distributed across \num{37} reporting authorities working to a common
directive, which is what makes them a property of how the measurement is
recorded rather than of who took it.

\begin{figure}[htbp]\centering
\includegraphics[width=0.86\textwidth]{fig04_loq_vs_eqs.pdf}
\caption{Substances for which European monitoring, as reported, cannot decide
the standard that applies to them: the share of station-years whose
quantification limit exceeds \num{30}\,\% of the environmental quality
standard, for substances with at least \num{1000} assessable station-years.
Generated by \texttt{scripts/90\_figures.py}.}
\label{fig:loqeqs}
\end{figure}

\subsection{The ontology applied to open data}
\label{sec:results-wbabox}
A vocabulary can be made to fit the dataset it was designed around; that proves
little. The test is whether its distinctions survive contact with a schema
written by someone else. Expressing a Waterbase subset in
CENSO\src{scripts/23\_waterbase\_abox.py} is that test: every class here is
assigned from EEA columns, and because the source is public the graph can be
rebuilt by any reader.

Of the \num{759257} river station-years whose substance carries a European
standard, a reservoir sample of \num{40000} was expressed in CENSO
(\num{59} analytes, \num{8695} stations, \num{441265} triples). Sampling rather
than inflating the graph to four million rows is deliberate: what is reasoned
over should be what is reported. The detection states partition cleanly --
\num{77.9}\% \texttt{CensoredObservation}, \num{10.0}\%
\texttt{QuantifiedObservation}, \num{12.2}\% \texttt{UnresolvedObservation} --
with the last of these the state no other surveyed vocabulary can express.

The compliance outcomes are the point. \num{62.1}\% are \texttt{Compliant} and
\num{2.3}\% \texttt{Exceedance}; the remaining \textbf{\num{35.7}\% are not
supportable}, split between \num{23.6}\% where the method's limit exceeds the
standard (\texttt{MethodInsufficient}, which Article~3(3b) requires to be set
aside) and \num{12.1}\% where the record carries no bound at all.

\paragraph{What substitution does to those \num{14268} assessments.}
It is tempting to describe the alternative as merely optimistic -- a pipeline
that calls everything it cannot decide \emph{compliant}. We assumed exactly
that in an earlier draft, and counting it showed the assumption to be wrong in
the direction that matters. Entering a non-detection at half its quantification
limit, the most widely used convention, pushes results \emph{above} the
standard: of the \num{10065} exceedances such a pipeline reports over these
\num{40000} records, only \num{908} rest on a quantified measurement above the
standard. The remaining \num{91}\,\% are artefacts of the substitution.

Nor is the choice of convention innocuous. Entering non-detections at zero
yields \num{3382} exceedances; at half the limit, \num{10065}; at the limit
itself, \num{12042}. The same records, the same standards and the same
arithmetic produce a \num{3.6}-fold range in the enforcement-relevant number,
selected by a rule the data does not determine (Figure~\ref{fig:verdicts}).
Reporting an interval and an explicit \emph{cannot be determined} is not a
refusal to answer; it is the only answer that does not depend on that choice.

\begin{figure}[htbp]\centering
\includegraphics[width=0.86\textwidth]{fig05_verdicts.pdf}
\caption{The same \num{40000} assessments, read four ways. The top three rows
are the identical records under the three substitution conventions in routine
use; they disagree with one another about the number of exceedances by a factor
of \num{3.6}, and only \num{908} of those exceedances rest on a quantified
measurement. The bottom row reports the two outcomes the law distinguishes
instead of absorbing them into compliance. Counts are computed row by row
rather than assumed; the plotted values are in
\texttt{supplementary/figure\_data/fig05\_verdicts.csv}. Generated by
\texttt{scripts/90\_figures.py}.}
\label{fig:verdicts}
\end{figure}

\paragraph{The two routes agree.}
Section~\ref{sec:results-waterbase} counts these failures with integer counters
and no ontology; this section reaches them by classification over a graph. The
method-insufficient share is \num{23.6}\% by both routes, against a sampling
standard error of \num{0.21} percentage points. That agreement is what makes
the vocabulary a \emph{representation} of the problem rather than a second
implementation of the same arithmetic: the counting could always have been
done, but only the graph can be queried for which assessments are unsupportable
and why.

\subsection{The same observations under two regulations}
\label{sec:results-dual}
Section~\ref{sec:methods-why} claims that pluggable regulation packages
let one set of observations be assessed under several jurisdictions without
reimplementation. The claim is cheap to make, so it is
run\src{scripts/24\_dual\_regulation.py}: both published packages are read as
data, and the same Waterbase rows are streamed past each. The observations, the
unit conversion and the decision procedure are identical in the two passes. The
only input that changes is which file supplies the threshold, and no code
changes at all.

Of \num{4190833} river station-years, \num{1551766} are regulated by at least
one of the two instruments. Reported naively, \num{82.8}\% of those change
outcome when the package is swapped -- but most of that is coverage: one
jurisdiction regulating a substance the other does not, which is a fact about
the two legal texts rather than about any measurement. The informative stratum
is the \num{325011} assessments \emph{both} jurisdictions regulate, where
nothing differs but the numeric standard. There, \textbf{\num{18.5}\%
(95\,\% CI \num{18.4}--\num{18.7}\,\%) still change outcome}: \num{51181} that
the European standard sets aside as method-insufficient are decidable under the
looser Turkish limit, \num{5372} move between \emph{exceeding} and
\emph{compliant} outright, and \num{3655} run the other way. Heavy metals and
the triazine herbicides dominate the divergence.

Two properties of that result matter more than its size. First, an assessment
that is \emph{method-insufficient} under one instrument and decidable under
another is not expressible at all if the threshold is a column: method adequacy
is a relation between the limit and the standard, so it changes when the
standard does. Second, the two verdicts can be held simultaneously, because the
thresholds carry different IRIs and different \texttt{definedBy} links --- the
graph records that the same measurement is compliant under one regime and
indeterminate under the other, rather than overwriting one with the other. Turkish
standards are of course not in force over European rivers; the comparison is
counterfactual by construction and is reported as a demonstration that the
machinery is regulation-independent, never as an assessment of Türkiye.

\subsection{Is the between-country spread an artefact?}
\label{sec:results-confounders}
Figure~\ref{fig:country-practice} shows the share of station-years carrying neither a
censoring flag nor a limit ranging from \num{3.3}\% to \num{97.4}\% between
countries reporting under the same directive. The obvious objection is that the
comparison is unfair: countries monitor different substances, over different
years, and some are not Member States at all. The objection is testable, so it
is tested\src{scripts/25\_country\_confounders.py} rather than argued with.

First, what is plotted is whether a schema field is populated. That field is
defined identically for every reporter by WISE-6; a national standard can change
which results breach a limit, but not whether the limit was written down.
Second, the spread survives every adjustment we can make to it. Standardising
each country directly onto the pooled European substance mix --- so that no
country can appear tidy merely because it happens to monitor substances usually
reported with a limit --- gives \num{8.9}--\num{97.3}\%. Restricting to the
\num{76} substances at least \num{60}\,\% of these countries report gives
\num{12.9}--\num{97.3}\%. Restricting to EU Member States, the only reporters
the quantification-limit criterion binds, gives \num{8.9}--\num{97.3}\%.
Kendall's~$\tau$ between the raw and standardised orderings is \num{0.89}: the
adjustment narrows the low end but does not reorder the countries. A spread of
that size, stable under three independent corrections, is a property of
reporting practice.

A fourth adjustment we intended --- restricting to recent years --- is not
available, and saying so matters more than the correction would have. Only
\num{8} of \num{37} reporting countries have any river record in this release
dated \num{2015} or later; France, Italy, Finland, Poland, Austria and Hungary
have none at all. The audit therefore describes the archived record, which for
most reporters ends before \num{2015}, and cannot speak to whether current
practice has improved. Whether it has is exactly the question a censoring-aware
schema would make answerable.

\subsection{Comparison with published ontologies}
\label{sec:results-gap}
% generated by scripts/07_verify_gap_table.py
% >>> begin tables/tab_gap.tex
% GENERATED by scripts/07_verify_gap_table.py -- do not edit.
\begin{table*}[htbp]\centering
\caption{Concepts present in published water and observation ontologies, decided by parsing each ontology file rather than its documentation. A dash means the file parsed and the concept was absent; \texttt{?} means it could not be retrieved, and no claim is made. The \emph{LOD bound to} column is the discriminating one: a detection limit declared as a sensor capability cannot state that a given measurement fell below it.}
\label{tab:gap}
\footnotesize
\begin{tabular}{l r l cccccc}
\toprule
Ontology & Ent. & LOD bound to & det.\ limit & censoring & undecid. & threshold & interval & applic. \\
\midrule
SOSA/SSN & 37 & -- & -- & -- & -- & -- & -- & -- \\
SOSA core & 53 & -- & -- & -- & -- & -- & -- & -- \\
GeoSPARQL & 0 & -- & -- & -- & -- & -- & -- & -- \\
QUDT schema & 239 & -- & -- & -- & -- & \checkmark & -- & \checkmark \\
SAREF core & 131 & -- & -- & -- & -- & \checkmark & -- & -- \\
SAREF4WATR v2.1.1 & 95 & -- & -- & -- & -- & \checkmark & -- & -- \\
SSN-system & 44 & the sensor & \checkmark & -- & -- & -- & -- & -- \\
OWL-Time & 78 & -- & -- & -- & -- & -- & -- & -- \\
WHOW water-monitoring & 57 & -- & -- & -- & -- & -- & \checkmark & -- \\
DoCE (Rio Doce WQ) & 168 & -- & -- & -- & -- & -- & -- & \checkmark \\
InWaterSense core & 189 & -- & -- & -- & -- & \checkmark & \checkmark & -- \\
InWaterSense regulations & 144 & -- & -- & -- & -- & -- & -- & -- \\
InWaterSense pollutants & 9 & -- & -- & -- & -- & -- & -- & -- \\
WAM-ONTO & 921 & -- & -- & -- & -- & -- & -- & -- \\
SewerNet & 145 & -- & -- & -- & -- & -- & -- & -- \\
ENVO & 7208 & -- & -- & -- & -- & -- & -- & -- \\
ExO (exposure) & 195 & -- & -- & -- & -- & -- & -- & -- \\
\textbf{CENSO (this work)} & 64 & \textbf{the result} & \checkmark & \checkmark & \checkmark & \checkmark & \checkmark & \checkmark \\
\textbf{CENSO-REG (this work)} & 28 & -- & -- & -- & -- & \checkmark & -- & -- \\
Project ontology 2018 & 123 & the sensor & \checkmark & -- & -- & -- & -- & -- \\
\bottomrule
\end{tabular}
\end{table*}
% <<< end tables/tab_gap.tex
\num{18} ontology files were assessed by parsing each one directly rather than
reading its documentation\src{scripts/07\_verify\_gap\_table.py}: \num{17}
published vocabularies plus one unpublished domain ontology, included so that a
non-standard baseline is held to the same test. Results are in
Table~\ref{tab:gap}. No ontology except CENSO carries a censoring status or an
outcome that records a verdict as not determinable.

Two of the remaining columns need reading carefully, because a keyword match is
not the same as a modelling commitment. Two vocabularies score for an
interval-valued result, but only WHOW attaches the interval to the
\emph{result}: InWaterSense's \texttt{hasRangeMinValue} and
\texttt{hasRangeMaxValue}, inherited from the SSN meteorology extension,
describe the range a \emph{device} can report. Likewise, among published
vocabularies only SSN-system defines a detection limit, and it binds it to the
sensor as a \texttt{SystemProperty}; the unpublished baseline does the same. In
both cases the term exists and the binding is to the instrument, which is why a
measurement cannot be said to have fallen below it.

A second dimension explains the pattern. Profiling each vocabulary for sampling
versus sensing terminology separates them cleanly: on a sampling/sensing count,
SAREF core scores \num{0}/\num{22} and InWaterSense core \num{2}/\num{27},
while SOSA core (\num{12}/\num{10}), WHOW (\num{10}/\num{0}) and CENSO
(\num{6}/\num{0}) lean the other way. \textbf{The field is dominated by
vocabularies built for in-situ sensing}, in which a detection limit is properly
a device property, and the few sampling-aware ones do not carry limits at all.

Five further ontologies described only in papers (WaWO+, OPO, InWaterSense's
expert system, Ahmedi et al., Acquirium) were assessed from their full
text\src{scripts/11\_assess\_competitor\_papers.py}: across all five, the terms
\emph{detection limit}, \emph{LOD}, \emph{LOQ}, \emph{censoring} and
\emph{non-detection} do not occur. This is weaker evidence than parsing a file
and is labelled as such.

\subsection{The ontology answers the questions it was built for}
\label{sec:results-cq}
Fifteen competency questions, each tied to an axiom, were executed over the
Waterbase graph together with the two vocabulary modules and both regulation
packages -- \num{442363} triples as
loaded\src{scripts/17\_run\_competency\_questions.py}. \textbf{All fifteen
return answers} (\num{25.4}\,s total). Because full OWL~2~RL closure over a
graph this size is impractical in pure Python, only two entailments are
materialised and both are stated in the report: union membership for the two
defined observation classes, and \texttt{rdfs:subClassOf} on type assertions.
Every other answer comes from asserted triples.

The questions are answered on public data, which matters more than the count:
a reader can re-run them without access to any private survey. Their SPARQL and
the rows returned are in the supplementary material.

\subsection{The axioms catch real defects}
Ten test cases, each injecting a defect present in this dataset, are all
rejected as intended\src{scripts/test\_axioms.py}: a record marked both
non-detect and quantified, a non-detection reported as an exceedance, a
half-LOQ substitution, one row joined to two substances by name matching, a
duplicated row, and a grab sample compared against an annual-average standard.
Nine are caught by OWL~2~RL entailment and one, outside the profile, by SHACL.

An OOPS!\ scan of the published distribution reports no critical pitfalls. Three
important ones were genuine and are fixed: properties lacking a domain or range.
The remainder are artefacts of packaging the modules as a single standalone file
for the scanner, which strips the imports and so leaves borrowed terms
undeclared; the distribution now re-declares them. Of the thirty-one properties
flagged as lacking an inverse, one genuinely did
(\texttt{definesThreshold}/\texttt{definedBy}) and is now declared; no inverse
was invented for the other thirty to satisfy the scanner.
% <<< end sections/05-results.tex
% >>> begin sections/06-discussion.tex
\section{Discussion}

\subsection{Relation to work on preprocessing sensitivity}
\label{sec:disc-lipinski}
Lipi\'nski \cite{lipinski2026sensitivity} recently examined the same underlying
problem from the analysis side, using \num{1259636} Waterbase observations of
cadmium, lead and nickel. Below-LOQ records accounted for \num{66.6}\% of
cadmium, \num{57.3}\% of lead and \num{36.1}\% of nickel observations, and
comparing three treatments -- discarding them, substituting half the LOQ, and
substituting the full LOQ -- substantially changed the inferred pollution
patterns. The recommendation is that studies report the rule they used.

That result is independent confirmation of our premise, on a different
continent-scale dataset, and we build on it rather than against it. Three
differences define what this paper adds.

\paragraph{The choice is treated as unavoidable; we remove it.}
All three of the scenarios compared yield a number. None can express that for a
given substance and threshold the data support no verdict at all. Our fourth
outcome is not another substitution rule but the absence of one, and it is
queryable: a regulator can ask which assessments are not supportable, and
receive a list.

\paragraph{The analysis presupposes that censoring survived to the record.}
Lipi\'nski could compare treatments because the rows examined carry a below-LOQ
flag and a \texttt{procedureLOQValue}. Our finding is that for a quarter of the
record they do not: \num{1066797} station-years (\num{25.5}\%) carry neither
(Section~\ref{sec:results-waterbase}). No sensitivity analysis can be run on
those rows, because there is nothing to be sensitive to --- not a choice between
substitution rules but the absence of anything to substitute from. The same gap
forced that study to \emph{exclude} below-LOQ records lacking an LOQ value.
Where it excludes them, CENSO represents them, as
\texttt{UnresolvedObservation}, and reports how many there are.

\paragraph{Metals are the mild case.}
The three substances studied are metals, and metals are where censoring bites
least. Across Waterbase, \num{47.1}\% of the \num{1491419} samples of
Annex~I metals fall below the quantification limit, against \num{77.3}\% of the
\num{4450281} samples of organic
micropollutants\src{scripts/22\_waterbase\_external.py}. The substance class
that recent European regulation has brought into scope is therefore censored at
close to twice the rate of the one the published sensitivity work examined, and
no comparable study has covered it.

\subsection{Failures that no preprocessing rule can repair}
Three of our findings are not preprocessing choices at all, and would survive
any substitution rule.

First, where the quantification limit exceeds the threshold the method cannot
demonstrate compliance however the non-detects are handled. This is the case for
\num{23.6}\% of all assessments against a European standard, for
\num{75.8}\% of imidacloprid assessments and for \num{100}\% of those for four
pyrethroids (Section~\ref{sec:results-waterbase}). Substituting zero, half the
limit or the full limit changes nothing: the measurement was never capable of
answering the question, and Section~\ref{sec:results-wbabox} shows what the
substitution does instead --- it manufactures a verdict whose direction depends
on the rule chosen.

Second, several thresholds carry preconditions the record cannot satisfy.
Cadmium's Annex~I standard is stated over five water-hardness classes, and the
bioavailability-corrected standards for nickel and lead require covariates that
the aggregated release does not carry alongside the measurement. Assessing such
a substance requires either the covariate or an explicit statement that the
standard could not be applied; a numeric threshold column offers neither, which
is why CENSO attaches preconditions to the threshold and returns
\texttt{PreconditionUnmet} rather than a default pass.

Third, a threshold can simply be transcribed wrongly, and nothing in a
conventional workflow would surface it. Both regulation packages released here
were transcribed from the primary legal texts and carry a machine-readable
\texttt{transcriptionStatus}; the parser that produced the European package is
verified against the consolidated
text\src{scripts/08\_verify\_thresholds.py}. That a threshold needs provenance
and a verification status at all is not a hypothetical concern: it is the same
argument as for the measurement, applied to the other side of the comparison.

\subsection{Why the vocabularies lacked this}
The gap is not an oversight. SOSA/SSN was built for in-situ sensing, where a
detection limit genuinely is a property of the device, and
\texttt{ssn-system:DetectionLimit} is correctly a \texttt{SystemProperty}. Trace
organic monitoring is a different modality: a sample is taken, transported and
analysed in a run with its own calibration and its own limits. The water
ontology literature inherited the sensing model, and our profiling shows the
inheritance directly -- the vocabularies that dominate the field score heavily
on sensing terminology and near zero on sampling. CENSO's contribution is to
carry the epistemics on the analytical run, where they belong.

\subsection{Limitations}
\begin{itemize}
  \item \textbf{We audit what was reported, not what was measured.} A missing
    limit may be a reporting omission rather than a laboratory one. That
    distinction matters for blame and not at all for the argument: either way
    the record cannot support the verdict drawn from it.
  \item \textbf{The graph is a sample.} Reasoning is over \num{40000}
    station-years drawn by reservoir sampling, not over all \num{4190833}.
    Inflating the graph to make a point about scale would misreport what was
    actually reasoned over; the population counts are given beside the sample
    throughout, and the two agree to within the sampling error.
  \item \textbf{Group standards are excluded from per-substance comparison.}
    A sum standard is not a limit on any of its members, so those rows are
    reported as not individually assessable rather than assessed.
  \item \textbf{One of the four outcomes is not exercised by this dataset.}
    \texttt{PossibleExceedance} --- the interval straddles the threshold, so
    the measurement cannot decide --- is instantiated \num{0} times in the
    Waterbase graph, and we would rather say so than let the four-valued claim
    rest on three values. There are two reasons, and neither is a defect of the
    vocabulary. For a \emph{quantified} result the class requires an
    uncertainty interval around the reported value, and the aggregated
    Waterbase release carries no per-result uncertainty: the interval is a
    point, so it cannot straddle anything. For a \emph{censored} result whose
    limit exceeds the standard the interval does straddle it, and both
    \texttt{PossibleExceedance} and \texttt{MethodInsufficient} apply --- but
    these answer different questions, the first epistemic and the second
    regulatory, and Article~3(3b) makes the regulatory answer the operative
    one. We report the outcome the law attaches, which is why those
    \num{9425} assessments appear under \texttt{MethodInsufficient}. A dataset
    reporting measurement uncertainty, as the disaggregated release does for
    some determinands, would exercise the fourth value directly; we did not use
    it here because an annual-average standard is defined against an annual
    mean, which is one aggregated row.
  \item \textbf{One matrix.} River water only; lakes, transitional and coastal
    waters have their own standards and are not analysed here.
  \item \textbf{The vocabulary comparison is of published files.} An ontology
    may be used in ways its published serialisation does not show, and three
    files could not be retrieved at all. Those are reported as unknown rather
    than as absent.
\end{itemize}

\subsection{Future work}
\label{sec:discussion-future}

\paragraph{Beyond compliance: attribution.}
The representation supports an attribution chain this paper does not pursue: a
substance's use category, taken from an external classification, mapped to an
expected pressure type and compared against the land use draining to a station.
What a knowledge-driven formulation would add over established source
apportionment -- positive matrix factorisation and chemical mass balance --
is not the idea but two properties: attribution that is \emph{explainable},
each assignment deriving from declared substance--use--sector relations rather
than from a factor interpreted after the fact, and attribution that is
\emph{censoring-aware}, able to return ``cannot be determined'' where
factorisation over
zero-substituted data returns a number. That comparison needs a denser station
network than any single survey here provides.

\paragraph{Ontology evaluation.}
All \num{15} competency questions were executed against the populated graph
with their SPARQL recorded, and all return non-empty
answers\src{scripts/17\_run\_competency\_questions.py}. The OOPS!\ pitfall
scanner reports no critical pitfall on the distributed
ontology\src{scripts/16\_export\_for\_scanners.py}; the one critical finding it
did raise -- an inverse pair whose domains made it incoherent -- was a genuine
modelling error and is fixed. A FOOPS!\ FAIR assessment has not yet been run,
and we do not report a FAIR score we have not obtained.

\subsection{Implications}
For monitoring programmes, the practical recommendation is narrow and cheap:
record the detection limit alongside every result, and record it per analytical
run. That single change makes the difference between a dataset on which
censoring can be reconstructed and one on which it cannot, and it costs nothing
at the point of measurement.

For regulators, the finding that an entire newly regulated substance class
cannot be assessed with the methods in routine use is actionable: an
environmental quality standard below the achievable limit of quantification is
unenforceable, and the mismatch is discoverable automatically once both are
machine-readable.

For the semantic web community, the case illustrates that a standard's implicit
measurement modality propagates into every extension built on it. SOSA's
sampling vocabulary exists but is not connected to analytical provenance, and
until it is, sample-and-analyse domains will keep reaching for a sensor model
that does not fit.
% <<< end sections/06-discussion.tex
% >>> begin sections/07-conclusions.tex
\section{Conclusions}

A non-detection is an upper bound and an open-world statement. Recording it as a
zero replaces that statement with a number, and the data models in routine use
offer nowhere else to put it.

In a four-campaign survey of \num{75} stations and \num{258} parameters,
\num{76.9}\% of chemical measurements -- \num{85.2}\% of micropollutant
measurements -- are stored as exactly zero, and the limits needed to reconstruct
what they mean exist for fewer than a third of the analytes. Beyond storage, we
found substances whose analytical method cannot demonstrate compliance in
principle, including every measured neonicotinoid under the European standard
now in force; thresholds whose preconditions the survey cannot satisfy; and a
working threshold table in which seven of ten metal standards diverged from the
law.

None of these states can be expressed in a model that admits only
\emph{compliant} and \emph{exceeding}. CENSO supplies the missing ones by
binding the detection limit to the analytical run that produced a result,
recording why a result is an interval, treating a threshold as a conditional
judgement with explicit preconditions, and making \emph{cannot be determined} a
first-class outcome. Its axioms were tested against ten defects drawn from this
dataset and reject all of them. A comparison of \num{18} published ontologies,
each parsed from its own file, shows that none combines these capabilities, and
that the gap follows from a measurement model built for sensors rather than for
samples.

The corrective for practice is small: record the detection limit with every
result, per analytical run. The corrective for representation is to stop
treating ``not detected'' as a number.
% <<< end sections/07-conclusions.tex

\section*{Data and code availability}
The ontology, the processing pipeline, every query and all evaluation outputs
are archived at \pending{Zenodo DOI}. The ontology and both regulation
packages are published under \url{https://w3id.org/censo/} (CC~BY~4.0); code
under MIT. The measurements are the EEA's Waterbase --- Water Quality ICM,
aggregated release, which is openly downloadable, so every figure and every
number in this paper can be re-derived without applying to anyone for
access.

\section*{CRediT author statement}
% DRAFT --- every role below must be confirmed by the named author before
% submission. Misattributing a CRediT role is a research-integrity matter, not
% a formatting one.
\textbf{Recep Can Altınbağ:} Conceptualization, Methodology, Software,
Formal analysis, Data curation, Validation, Visualization, Writing --- original
draft. \textbf{Mehmet Tahir Sandıkkaya:} Methodology, Supervision,
Writing --- review \& editing. \textbf{Ulaş Pekyavaş:} Investigation,
Writing --- review \& editing.

\section*{Declaration of generative AI and AI-assisted technologies
in the writing process}
% Elsevier requires this statement when a generative tool assisted the WRITING.
% Fill in the tool and the purpose truthfully; the authors remain responsible
% for every claim regardless of what assisted its drafting.
During the preparation of this work the authors used
\pending{tool and purpose}. After using this tool the authors reviewed and
edited the content as needed and take full responsibility for the content of
the publication.

No figure in this article was generated or altered by a generative model. Every
figure is drawn by \texttt{scripts/90\_figures.py} from the tabular data
deposited with the article, the values behind each mark are shipped in
\texttt{figure\_data/}, and any reader can regenerate the figures from the
public source data without contacting the authors.

\bibliographystyle{unsrt}
\bibliography{refs}

\end{document}
